\documentclass[../DoAn.tex]{subfiles}
\begin{document}

\begin{center}
    \Large{\textbf{ABSTRACT}}\\
\end{center}
\vspace{1cm}

In modern administrative management, automated access control and attendance systems face significant challenges regarding identity fraud risks, the logistical burden of issuing internal access cards, and the risk of biometric data leakage when storing portrait image repositories on central servers. To comprehensively address these limitations, this thesis proposes and implements a distributed, autonomous two-layer authentication system that directly leverages Vietnam's existing national identity infrastructure — the chip-embedded Citizen Identity Card (CCCD). This approach allows the system to legally extract the original DG2 biometric image partition from the card's security chip via NFC near-field communication at the moment of scanning, using it as a real-time reference template for comparison against a live selfie, thereby completely eliminating the cumbersome manual sample-collection phase.

The proposed system architecture comprises five asynchronously operating core components: an Android Flutter terminal application that performs NFC chip reading and runs a compact MobileFaceNet model for first-layer on-device verification; a central Java Spring Boot backend engine that handles access permission matrix validation; a Python Flask AI server that performs feature vector extraction and cosine similarity computation using the deep convolutional ArcFace network for the second, in-depth authentication layer; an ESP32 embedded microcontroller that receives JSON commands via the lightweight MQTT protocol to actuate a relay and release the physical door lock; and a React Web Dashboard that automatically updates a real-time monitoring log through a full-duplex WebSocket STOMP pipeline.

The principal contribution of this thesis is the successful implementation and proof-of-concept demonstration of a stateless dual-layer attendance architecture that stores no centralized portrait image repository, thereby protecting the personal data privacy of end users. Experimental results confirm that the system achieves a high level of readiness, operating stably with an average end-to-end response latency of 3.3 seconds, and is ready for real-world deployment in organizations and enterprises.

\end{document}